% LSPV-5 NOTE::DIGCORR1::2026-10-02T07:58:00-07:00
\documentclass[11pt]{article}\usepackage[margin=.88in]{geometry}\usepackage{amsmath,amssymb}
\begin{document}
\section*{DIGCORR1: digit-correlation route and multiplicative-shift exceptional set}
\textbf{STATUS: exact algebraic sublemma + literature-supported candidate analytic route; P versus NP remains OPEN.}
\paragraph{Literature synthesis.}
Correlation-measure work studies dependence among shifted binary sequences, including finite-field polynomial/character constructions. Digital-number-theory work shows non-algebraic digit functions can sometimes be controlled by exponential sums. Bit-security work represents standard bit predicates modulo odd moduli through segment predicates/Fourier analysis. Multiplicative-subgroup literature gives nontrivial additive-character cancellation. None directly proves the required rank theorem for \(b_j(t)=\operatorname{bit}_j(f(g^t))\).
\paragraph{Rank target.}
BILINREC1 gives \(d=O(m)\Rightarrow\dim\ker\Phi=\Theta(m^2)\). It suffices instead to exhibit \(\omega(m)\) independent rows/constraints. A possible route is to select many bit-pair/lag correlation streams whose signed lifts are pairwise almost orthogonal; Gram-matrix rank bounds would then force \(\operatorname{rank}\Phi=\omega(m)\).
\paragraph{Fourier reduction candidate.}
Each standard bit predicate on representatives \(0,\ldots,p-1\) is a union of intervals (with complexity depending on bit position). Fourier expansion of interval indicators reduces shifted bit correlations to additive-character sums whose phases include
\[
F_{\alpha,\beta,\delta}(x)=\alpha f(x)+\beta f(g^\delta x).
\]
Classical Weil-type bounds are potentially strong because \(\deg f<h\) while vMS takes \(p\) enormous relative to \(h\).
\paragraph{PROVED/EXACT exceptional-set criterion.}
Write \(f(x)=\sum_{e\in E}c_ex^e\), with \(c_e\ne0\). Then
\[
F_{\alpha,\beta,\delta}(x)=\sum_{e\in E}c_e(\alpha+\beta g^{\delta e})x^e.
\]
For all nonconstant terms to cancel with \(\alpha\beta\ne0\), every nonzero active exponent \(e\) must obey the same equality \(g^{\delta e}=-\alpha/\beta\). Hence for all active nonzero exponents \(e,e'\),
\[
\boxed{g^{\delta(e-e')}=1}
\]
or equivalently
\[
\boxed{(p-1)\mid\delta(e-e')}.
\]
Thus multiplicative-shift Fourier degeneracy is confined to explicit congruence classes determined by exponent differences. Since all exponents are below \(h\ll p\), sufficiently small nonzero lags avoid degeneracy whenever \(f\) has two nonzero nonconstant exponents with difference not annihilated modulo \(p-1\).
\paragraph{Exceptional sparse case.}
If \(f\) has only one nonconstant monomial (plus possibly a constant), the preceding two-exponent obstruction is unavailable. This sparse-support class must be handled separately. Boolean-grid constraints may strongly restrict such polynomials, but no exclusion theorem is claimed here.
\paragraph{Open analytic bridge.}
Even for nondegenerate \(F_{\alpha,\beta,\delta}\), converting additive-character cancellation into sufficiently uniform bounds for all Fourier terms of individual bit predicates and then into \(\mathbb F_2\)-rank of the supported operator remains OPEN. Fourier truncation cost, high/low bit segment complexity, exceptional frequencies, repeated-root support, and same-shift legality must all be charged.
\paragraph{Next target.}
First classify whether a nonconstant Boolean-grid polynomial \(f|_H\in\{0,1\}\), \(\deg f<h\), can be monomial-plus-constant when \(h\) is large. If this sparse exceptional class is impossible or tiny, then quantify the number of small lags \(\delta<m\) for which every nontrivial two-frequency phase remains nonconstant, and test the resulting correlation matrix numerically before attempting the full exponential-sum/rank proof.
\end{document}
% END_OF_VERIFIED_NOTE::LSPV-5::DIGCORR1